\documentclass[11pt]{article}

\usepackage[margin=1in]{geometry}
\usepackage{amsmath,amssymb}
\usepackage{booktabs}
\usepackage{enumitem}
\usepackage{hyperref}
\hypersetup{hidelinks}

\title{Scaling Book Exercises: Section 12, Quiz 3}
\author{Lucas Yan}
\date{}

\begin{document}
\maketitle

\noindent
\textbf{Quiz:} Beyond the Node Level\\
\textbf{Source scan:} \texttt{../handwritten/section\_12\_handwritten.pdf},
pages 5--6\\
\textbf{Reference:} \url{https://jax-ml.github.io/scaling-book/gpus/}

\section*{Question 1: Fat-tree topology}

The reference H100 SuperPod has \(1024\) GPUs, or \(128\) eight-GPU nodes.
Each node connects to the leaf layer using eight \(400\) Gbit/s links:
\[
8\cdot\frac{400}{8}
=\boxed{400\ \text{GB/s per node}}.
\]

A scalable unit (SU) contains \(32\) nodes and eight leaf switches.  The
node-facing half of the \(64\) ports on each leaf switch provides
\[
8\cdot32\cdot\frac{400}{8}
=\boxed{12.8\ \text{TB/s per SU}},
\]
again \(400\) GB/s for each of its 32 nodes.

The leaf-to-spine links provide the same SU bandwidth:
\[
8\text{ leaf}\cdot16\text{ spine}\cdot2\text{ links}
\cdot\frac{400}{8}
=\boxed{12.8\ \text{TB/s per SU}}.
\]
Finally, the 16 spine switches each support
\[
64\cdot\frac{400}{8}=3.2\ \text{TB/s},
\]
so their aggregate capacity is
\[
16(3.2)=\boxed{51.2\ \text{TB/s}}.
\]

Bisecting the 128 nodes leaves 64 nodes on each side.  Counting both
directions,
\[
2\cdot64\cdot400\ \text{GB/s}
=\boxed{51.2\ \text{TB/s}}.
\]
Every level supplies \(400\) GB/s per node, so no level oversubscribes the
others: the network has full bisection bandwidth.

\section*{Question 2: Scaling the pod}

\paragraph{2048 GPUs.}
Keep the 32-node SU unchanged, increase from four to eight SUs, and double
the spine layer from 16 to 32 switches.  A leaf switch has only 64 ports,
so use one \(400\) Gbit/s link from each leaf to each spine rather than two;
doubling the number of spine switches preserves total bandwidth.

\paragraph{4096 GPUs.}
The leaf/spine design then runs out of ports.  Add another fat-tree level
of core switches.  One reference arrangement uses
\[
\boxed{128\text{ spine switches and }64\text{ core switches}.}
\]
The essential principle is to add switching capacity at the same rate as
node bandwidth so that the larger fabric remains non-oversubscribed.

\end{document}
